\chapter{Stochastic Quantization Reborn: AI Sampling}

The newest pillar of the library is machine-learning sampling. Its
physical root is forty years older: stochastic quantization represents
quantum field theory as Langevin diffusion, and generative diffusion
models are exactly time-reversed Langevin dynamics. Two original articles
of the library founded the lattice application.

% ----------------------------------------------------------------------
\section{Stochastic Quantization}
\papercard{G.\ Parisi, Y.-S.\ Wu}
{Perturbation theory without gauge fixing}
{Sci.\ Sinica \textbf{24}, 483 (1981)}
{No arXiv (preprint era)}
\whyselected{Root of the entire AI-sampling line ($\sim$3/99 within the
AI corpus, plus conceptual citations in flow papers): the idea that a
quantum field theory equals the equilibrium statistics of a stochastic
process predates, and now powers, the machine-learning era.}
\physics{Promote fields to functions of a fictitious fifth time $s$
governed by Langevin dynamics,
\begin{equation}
  \frac{\partial A_\mu(x,s)}{\partial s}
   = -\frac{\delta S[A]}{\delta A_\mu(x,s)} + \eta_\mu(x,s),
   \qquad
   \langle\eta_\mu(x,s)\eta_\nu(y,s')\rangle = 2\delta_{\mu\nu}\delta(x-y)\delta(s-s'),
\end{equation}
whose equilibrium distribution is $e^{-S}$. Because the drift is gauge
covariant, no gauge fixing is needed --- perturbation theory proceeds in
the physical (transverse) subspace automatically, and the Fokker--Planck
operator generates the renormalization group (Zinn--Justin). The flow-time
structure will reappear verbatim in L\"uscher's gradient flow, with
deterministic drift replacing the noise average.}
\impact{Wang--Aarts--Zhou's diffusion-as-stochastic-quantization paper
(in the library) makes the correspondence exact and constructive:
training a score-based diffusion model reproduces the reverse Langevin
process, yielding Boltzmann-distributed samples of the field theory.}
\cardend

% ----------------------------------------------------------------------
\section{Flow-Based MCMC}
\papercard{M.\ S.\ Albergo, G.\ Kanwar, P.\ E.\ Shanahan}
{Flow-based generative models for Markov chain Monte Carlo in lattice
field theory}
{Phys.\ Rev.\ D \textbf{100}, 034515 (2019)}
{arXiv:1904.12072 | DOI: 10.1103/PhysRevD.100.034515}
\whyselected{Founding article of machine-learning sampling for lattice
field theory, and an original article of this library ($\sim$4/99 within
its niche): cited by every AI-for-lattice paper of the corpus.}
\physics{Train an invertible neural network $f_\phi$ as a proposal
distribution $q_\phi(\phi)=p_0(f_\phi(\phi))|\det J_{f_\phi}|$, then use
it inside Metropolis--Hastings with acceptance
$\min(1, e^{-S[\phi]}/e^{-S[f_\phi(\phi)]}\cdot q/p)$: exactness of the
target distribution survives any imperfection of the learned model.
Demonstrated on 2d $\phi^4$ theory, the approach suppresses critical
slowing-down precisely where local algorithms fail, and trains without
samples of the target distribution --- only the action enters the loss.
This resolves the classic chicken-and-egg problem of simulation-based
learning.}
\impact{Within the library this paper is the common ancestor of the
equivariant-gauge extension below, of the diffusion-as-stochastic-
quantization line, and of the physics-driven-learning agenda for inverse
problems; conceptually it realizes the trivializing-map dream of
Chapter~3 with learnable maps.}
\cardend

% ----------------------------------------------------------------------
\section{Gauge Equivariance at Scale}
\papercard{G.\ Kanwar, M.\ S.\ Albergo, D.\ Boyda, K.\ Cranmer,
D.\ C.\ Hackett, S.\ Racani\`ere, D.\ J.\ Rezende, P.\ E.\ Shanahan}
{Equivariant flow-based sampling for lattice gauge theory}
{Phys.\ Rev.\ Lett.\ \textbf{125}, 121601 (2020)}
{arXiv:2003.06413 | DOI: 10.1103/PhysRevLett.125.121601}
\whyselected{Original article of this library; the demonstration that
machine learning can attack the hardest sampling problem of lattice QCD
--- topology ($\sim$3/99 within its niche).}
\physics{Build the flow from layers that commute with gauge
transformations (maskings coupling links to plaquette-derived functions,
coupling kernels projected onto the group through exponentiated traces),
so that $q_\phi(U^g)=q_\phi(U)$ and the sampled ensemble is manifestly
gauge invariant. Applied to 2d $U(1)$ pure gauge near the continuum
limit, trained models reduce the autocorrelation time of the topological
charge by factors of $\sim1500$ relative to HMC and $\sim200$ relative to
heat bath --- attacking the exponential freezing of topology that plagues
fine-lattice QCD simulations cited since L\"uscher's own work on open
boundary conditions.}
\impact{The library frames its AI-for-QCD narrative around this result:
symmetry-respecting architectures are not a constraint but the source of
the gain --- the same lesson its agentic-computing papers draw when they
insist that deterministic domain tools must remain in charge of the
physics while language models orchestrate. The circle back to
Parisi--Wu and to trivializing maps closes here.}
\cardend
